% \chapter*{Glossario}

% \paragraph{Open source} In informatica, con l'espressione open source si indica un software distribuito, generalmente in via gratuita, sotto i termini di una licenza open source, che ne concede lo studio, l'utilizzo, la modifica e la redistribuzione. Questo modello si pone in contrapposizione con l'idea di software proprietario che permette le sopracitate concessioni solo secondo i termini dettati dal detentore del copyright.

% \paragraph{JVM} La macchina virtuale Java (detta anche Java Virtual Machine o JVM) è il componente della piattaforma Java responsabile per l'esecuzione dei programmi in formato bytecode. La JVM viene utilizzata per eseguire codice Java.

% \newpage

\newglossaryentry{Middleware}{
    name={Middleware},
    description={Il middleware è un software che funge da strato intermedio tra le applicazioni e le componenti sottostanti, come ad esempio sistemi operativi, database o hardware, facilitando la comunicazione, l'integrazione e la gestione delle risorse nei sistemi distribuiti. Il suo ruolo primario è quello di nascondere la complessità dell'infrastruttura sottostante, garantendo interoperabilità, scalabilità e sicurezza}    
}

\newglossaryentry{Open Source}{
    name={Open Source},
    description={software distribuito, generalmente in via gratuita, sotto i termini di una licenza open source}
}

\newglossaryentry{Topic Kafka}{
    name={Topic Kafka},
    description={I Topic di Kafka sono le categorie utilizzate per organizzare i messaggi. Ogni topic ha un nome unico per l'intero cluster Kafka.
    I messaggi vengono inviati e letti da argomenti specifici. I produttori scrivono dati negli argomenti e i consumatori leggono dati dagli argomenti. 
    In Kafka, gli argomenti sono partizionati e replicati attraverso i broker per tutta l'implementazione.  I broker si riferiscono a ciascuno dei nodi di un cluster Kafka. Le partizioni sono importanti perché permettono la parallelizzazione degli argomenti, consentendo un elevato throughput dei messaggi}
}

\newglossaryentry{Visual Studio Code}{
    name={Visual Studio Code},
    description={editor di codice sorgente sviluppato da Microsoft per Windows, Linux e macOS. È un software libero e gratuito per uso personale e commerciale, anche se la versione ufficiale è sotto una licenza proprietaria. Permette di installare estensioni che aggiungono ulteriori funzionalità}
}
